% Provenance (never printed): see sources/outline_sources.md and sources/book1/BRIEFS.md (g2/02).
\chapter{Using a Computer}\label{ch:g2:using-a-computer}

Sami spent the whole afternoon drawing a castle on the \omterm{def:g1:what-is-a-computer:computer}{computer}: towers, a
bridge, a dragon on the roof. When his mother called him for dinner, he closed
the drawing and switched off the \omterm{def:g1:what-is-a-computer:computer}{computer}. The next morning the castle was
gone. The \omterm{def:g1:what-is-a-computer:computer}{computer} had not kept it, because Sami never asked it to.

To use a \omterm{def:g1:what-is-a-computer:computer}{computer} well, we need to know what we see on the screen, how to
start and stop our work, and how to make sure the \omterm{def:g1:what-is-a-computer:computer}{computer} keeps it.

\begin{recall}
The \omterm{def:g2:keyboard-and-mouse:pointer}{pointer} moves with the mouse; a \omterm{def:g2:keyboard-and-mouse:click}{click} chooses something, a \omterm{def:g2:keyboard-and-mouse:click}{double-click}
often opens it, and dragging moves it (\cref{ch:g2:keyboard-and-mouse}).
\end{recall}

\section{Switching on and shutting down}

To switch a \omterm{def:g1:what-is-a-computer:computer}{computer} on, press its power button once. The \omterm{def:g1:what-is-a-computer:computer}{computer} then needs
some time to get ready: the screen lights up, and after a little while the
\omterm{def:g2:using-a-computer:desktop}{desktop} appears. Wait until it is ready before you \omterm{def:g2:keyboard-and-mouse:click}{click}.

When you have finished, do not simply press the power button or pull out the
plug. Shut the \omterm{def:g1:what-is-a-computer:computer}{computer} down from its \omterm{def:g2:using-a-computer:menu}{menu}: it then finishes writing what it
was keeping, closes everything properly, and switches itself off.

\begin{safety}
Save your work often, not only at the end. Shut the \omterm{def:g1:what-is-a-computer:computer}{computer} down from its
\omterm{def:g2:using-a-computer:menu}{menu}, never by pulling out the plug: work that was being written at that moment
can be lost.
\end{safety}

\section{The desktop, icons and windows}

\begin{definition}[Desktop and icon]\label{def:g2:using-a-computer:desktop}
The \emph{desktop}\index{desktop} is the background of the screen when the
\omterm{def:g1:what-is-a-computer:computer}{computer} is ready, the place where we start our work. An
\emph{icon}\index{icon} is a small picture that stands for something we can
open: a piece of work or a tool.
\end{definition}

\omterm{def:g2:keyboard-and-mouse:click}{Double-click} on an \omterm{def:g2:using-a-computer:desktop}{icon}, and what it stands for opens in a \omterm{def:g2:using-a-computer:window}{window}.

\begin{definition}[Window]\label{def:g2:using-a-computer:window}
A \emph{window} is a rectangle on the screen that shows one piece of work or
one tool, with a title bar at the top.
\index{window}
\end{definition}

Several \omterm{def:g2:using-a-computer:window}{windows} can be open at the same time. They may overlap, like sheets
of paper on a desk: the \omterm{def:g2:using-a-computer:window}{window} in front is the one you are working in. \omterm{def:g2:keyboard-and-mouse:click}{Click}
on a \omterm{def:g2:using-a-computer:window}{window} that is behind, and it comes to the front.

\begin{omfigure}
\begin{tikzpicture}[x=1cm, y=1cm, font=\scriptsize]
  % the screen and desktop
  \draw[black!70, fill=ocPaleBlue!60, rounded corners=2pt] (0,0) rectangle (12,7);
  \foreach \y/\t/\c in {6.0/drawing/ocRed!55, 4.6/story/ocBlue!55, 3.2/music/ocMint!70} {
    \draw[black!60, fill=\c, rounded corners=2pt] (0.45,\y-0.25) rectangle (1.15,\y+0.45);
    \node at (0.8,\y-0.5) {\t};
  }
  % window behind
  \omwindow{6.2}{4.2}{Story}
  \pic at (2.0,1.9) {window};
  \foreach \y in {5.2,4.85,...,2.4} \draw[black!25, line width=1.2pt] (2.3,\y) -- (5.2,\y);
  % window in front
  \omwindow{6.0}{4.3}{Drawing}
  \pic at (5.4,0.6) {window};
  \draw[ocRed!70, line width=1pt] (6.3,1.3) rectangle (7.8,2.9);
  \draw[ocRed!70, line width=1pt] (6.3,2.9) -- (6.3,3.3) -- (6.6,3.3) -- (6.6,2.9);
  \draw[ocRed!70, line width=1pt] (7.5,2.9) -- (7.5,3.3) -- (7.8,3.3) -- (7.8,2.9);
  \draw[ocBlue!70, line width=1pt] (8.3,1.3) -- (10.6,1.3) -- (10.6,2.4) -- (9.45,3.3) -- (8.3,2.4) -- cycle;
  % labels outside the screen
  \node[anchor=west, font=\footnotesize] (li) at (12.3,6.1) {icons};
  \draw[black!60] (12.25,6.1) -- (1.2,6.1);
  \node[anchor=west, font=\footnotesize] at (12.3,3.7) {window in front};
  \draw[black!60] (12.25,3.7) -- (11.35,3.7);
  \node[anchor=west, font=\footnotesize] at (12.3,5.0) {window behind};
  \draw[black!60] (12.25,5.0) -- (8.1,5.75);
  \node[anchor=west, font=\footnotesize] at (12.3,0.3) {desktop};
  \draw[black!60] (12.25,0.3) -- (2.0,0.3);
\end{tikzpicture}

{\small A \omterm{def:g2:using-a-computer:desktop}{desktop} with three \omterm{def:g2:using-a-computer:desktop}{icons} and two open \omterm{def:g2:using-a-computer:window}{windows}. The \omterm{def:g2:using-a-computer:window}{window} in front,
``Drawing'', hides part of the \omterm{def:g2:using-a-computer:window}{window} ``Story''.}
\label{fig:g2:using-a-computer:desktop}
\end{omfigure}

Every \omterm{def:g2:using-a-computer:window}{window} has the same parts. You move a \omterm{def:g2:using-a-computer:window}{window} by dragging its title bar.
You make it bigger or smaller by dragging one of its corners. Three small
buttons at one end of the title bar hide the \omterm{def:g2:using-a-computer:window}{window} (it waits, out of the way,
until you call it back), make it fill the whole screen, or close it.

\begin{definition}[Menu]\label{def:g2:using-a-computer:menu}
A \emph{menu}\index{menu} is a list of choices that opens when you \omterm{def:g2:keyboard-and-mouse:click}{click} on
its name; clicking on a choice carries it out.
\end{definition}

We write a choice in a \omterm{def:g2:using-a-computer:menu}{menu} with an arrow: \menu{File > Save} means ``\omterm{def:g2:keyboard-and-mouse:click}{click} on
the \omterm{def:g2:using-a-computer:menu}{menu} \menu{File}, then on the choice \menu{Save}''.

\begin{omfigure}
\begin{tikzpicture}[x=1cm, y=1cm, font=\footnotesize, lab/.style={font=\footnotesize, anchor=west}]
  \omwindow{8}{5}{My castle}
  \pic at (0,0) {window};
  % menu bar
  \draw[black!40] (0,4.05) -- (8,4.05);
  \node[anchor=west] at (0.1,4.3) {File\quad Edit\quad View};
  % open File menu
  \draw[black!60, fill=white] (0.12,2.35) rectangle (1.9,4.05);
  \fill[ocYellow!50] (0.14,2.88) rectangle (1.88,3.22);
  \foreach \y/\t in {3.75/New, 3.4/Open, 3.05/Save, 2.7/Save as} \node[anchor=west] at (0.2,\y) {\t};
  % work area drawing
  \draw[ocRed!70, line width=1pt] (3.2,0.6) rectangle (5.6,2.6);
  \draw[ocRed!70, line width=1pt] (3.2,2.6) -- (3.2,3.1) -- (3.6,3.1) -- (3.6,2.6) (5.2,2.6) -- (5.2,3.1) -- (5.6,3.1) -- (5.6,2.6);
  % resize corner
  \foreach \d in {0.12,0.24,0.36} \draw[black!50] (8-\d,0.05) -- (7.95,\d);
  % labels
  \draw[black!60] (4,4.78) -- (4,5.6) node[above] {title bar};
  \draw[black!60] (7.45,4.78) -- (8.8,5.3) node[lab] {close};
  \draw[black!60] (6.9,4.78) -- (8.8,4.85) node[lab] {fill the screen};
  \draw[black!60] (6.35,4.78) -- (8.8,4.4) node[lab] {hide};
  \draw[black!60] (0.05,4.3) -- (-0.6,4.3) node[anchor=east] {menus};
  \draw[black!60] (1.0,3.05) -- (-0.6,3.05) node[anchor=east] {an open menu};
  \draw[black!60] (5.6,1.6) -- (8.8,2.2) node[lab] {the work};
  \draw[black!60] (7.9,0.15) -- (8.8,0.4) node[lab] {corner for resizing};
\end{tikzpicture}

{\small The parts of a \omterm{def:g2:using-a-computer:window}{window}. The \omterm{def:g2:using-a-computer:menu}{menu} \menu{File} is open, and the choice
\menu{Save} is under the \omterm{def:g2:keyboard-and-mouse:pointer}{pointer}.}
\label{fig:g2:using-a-computer:window-parts}
\end{omfigure}

\section{Starting a program and writing a text}

Each \omterm{def:g2:using-a-computer:desktop}{icon} of the \omterm{def:g2:using-a-computer:desktop}{desktop}, or each line of the main \omterm{def:g2:using-a-computer:menu}{menu} of the \omterm{def:g1:what-is-a-computer:computer}{computer},
starts a tool: a tool for drawing, one for writing, one for listening to
music. The tool opens in a \omterm{def:g2:using-a-computer:window}{window}, and closing the \omterm{def:g2:using-a-computer:window}{window} stops it.

To write a text, there are two kinds of tool.

\begin{definition}[Text editor and word processor]\label{def:g2:using-a-computer:word-processor}
A \emph{text editor}\index{text editor} is a tool for writing that keeps only
the letters, digits, signs, spaces and new lines of a text. A \emph{word
processor}\index{word processor} also keeps how the text looks: words in bold,
sizes and colours of letters, pictures, the layout of the page.
\end{definition}

\begin{example}[The same note in two tools]\label{ex:g2:using-a-computer:note}
Mia writes an invitation to her party. In a \omterm{def:g2:using-a-computer:word-processor}{text editor}, every letter looks
the same: the title cannot be made bigger. In a \omterm{def:g2:using-a-computer:word-processor}{word processor}, she makes the
title big and bold, colours the date red, and adds a picture of a cake. Many
\omterm{def:g2:using-a-computer:word-processor}{word processors} also underline words that look misspelt, so that she can
check them.
\end{example}

\begin{omfigure}
\begin{tikzpicture}[x=1cm, y=1cm, font=\scriptsize]
  \omwindow{6.6}{4.2}{Text editor}
  \pic at (0,0) {window};
  \node[anchor=north west, align=left, font=\ttfamily\scriptsize] at (0.2,3.6)
    {My party\\[2pt]Come to my party on\\Saturday at 3 o'clock.\\There will be cake!\\[2pt]Mia};
  \omwindow{6.6}{4.2}{Word processor}
  \pic at (7.2,0) {window};
  \node[anchor=north west, align=left] at (7.4,3.6)
    {{\Large\bfseries My party}\\[3pt]{\small Come to my party on}\\{\small\color{ocRed}\bfseries Saturday at 3 o'clock.}\\{\small There will be cake!}\\[2pt]{\small\itshape Mia}};
  \fill[ocYellow!70] (12.0,1.0) rectangle (13.2,1.6);
  \fill[ocRed!60] (12.0,1.6) rectangle (13.2,1.75);
  \foreach \x in {12.2,12.6,13.0} \draw[ocRed!80, line width=1pt] (\x,1.75) -- (\x,2.05);
\end{tikzpicture}

{\small The same invitation in a \omterm{def:g2:using-a-computer:word-processor}{text editor} (left), which keeps only the
letters, and in a \omterm{def:g2:using-a-computer:word-processor}{word processor} (right), which also keeps sizes, colours and
pictures.}
\label{fig:g2:using-a-computer:editor-processor}
\end{omfigure}

\section{Saving and finding work again}

While you work, what you do lives only inside the \omterm{def:g1:what-is-a-computer:computer}{computer} while it is
switched on. If the \omterm{def:g2:using-a-computer:window}{window} is closed or the \omterm{def:g1:what-is-a-computer:computer}{computer} switched off without
saving, the work is lost: this is what happened to Sami's castle.

\begin{definition}[Saving]\label{def:g2:using-a-computer:save}
\emph{Saving}\index{saving (a file)} a piece of work means asking the
\omterm{def:g1:what-is-a-computer:computer}{computer} to keep a copy of it, under a name and at a place we choose, so that
it is still there after the \omterm{def:g1:what-is-a-computer:computer}{computer} has been switched off.
\end{definition}

The first time, choose \menu{File > Save as}: the \omterm{def:g1:what-is-a-computer:computer}{computer} asks for a name and
for a place. After that, \menu{File > Save} keeps the newest version under the
same name, at the same place. To keep the old version and make a new copy,
choose \menu{File > Save as} again and give a new name.

\begin{method}[Choosing a good name]\label{met:g2:using-a-computer:name}
A good name is short and says what the work is: \code{castle drawing} rather
than \code{aaa} or \code{my work}. If you keep several versions, add a number or
a date: \code{story version 2}, \code{story 2026-03-14}.
\end{method}

\begin{omfigure}
\begin{tikzpicture}[x=1cm, y=1cm, font=\footnotesize]
  \omwindow{8.2}{4.4}{Save as}
  \pic at (0,0) {window};
  \node[anchor=west] at (0.3,3.45) {Name:};
  \draw[black!60, fill=white] (1.4,3.2) rectangle (7.8,3.7);
  \node[anchor=west, font=\ttfamily\footnotesize] at (1.5,3.45) {castle drawing};
  \node[anchor=west] at (0.3,2.75) {Place:};
  \draw[black!60, fill=white] (1.4,0.9) rectangle (4.6,2.95);
  \foreach \y/\t in {2.7/Desktop, 2.3/Documents, 1.9/Pictures, 1.5/School work} \node[anchor=west] at (1.5,\y) {\t};
  \fill[ocBlue!25] (1.42,1.7) rectangle (4.58,2.1); \node[anchor=west] at (1.5,1.9) {Pictures};
  \draw[black!60, fill=ocBlue!35, rounded corners=2pt] (6.4,0.3) rectangle (7.8,0.8);
  \node at (7.1,0.55) {Save};
  \draw[black!60, fill=white, rounded corners=2pt] (4.8,0.3) rectangle (6.2,0.8);
  \node at (5.5,0.55) {Cancel};
\end{tikzpicture}

{\small \omterm{def:g2:using-a-computer:save}{Saving} a drawing for the first time: a good name, and a place chosen in
the list. The work will be found there again.}
\label{fig:g2:using-a-computer:save-as}
\end{omfigure}

\begin{omfigure}
\includegraphics[width=0.55\linewidth]{images/book1/ai/g2-using-a-computer-saving.jpg}

{\small \omterm{def:g2:using-a-computer:save}{Saving} work at school, with a teacher's help.}
\end{omfigure}

To find your work again the next day, open the tool you used, choose
\menu{File > Open}, and look at the place where you saved it. Or look at that
place directly and \omterm{def:g2:keyboard-and-mouse:click}{double-click} on the name of your work. A good name makes it
easy to recognise.

\section{Exercises}

\begin{exercise}[$\star$]\label{exo:g2:using-a-computer:1}
Name four parts of a \omterm{def:g2:using-a-computer:window}{window}.
\end{exercise}

\begin{exercise}[$\star$]\label{exo:g2:using-a-computer:2}
How do you move a \omterm{def:g2:using-a-computer:window}{window} to another place on the screen? How do you make it
bigger?
\end{exercise}

\begin{exercise}[$\star$]\label{exo:g2:using-a-computer:3}
Why should we shut a \omterm{def:g1:what-is-a-computer:computer}{computer} down from its \omterm{def:g2:using-a-computer:menu}{menu} rather than pull out the
plug?
\end{exercise}

\begin{exercise}[$\star$]\label{exo:g2:using-a-computer:4}
Sort these names into good names and bad names for a piece of work:
\code{aaaa}, \code{volcano poster}, \code{work}, \code{class photo 2026},
\code{xyz2}, \code{letter to grandma}.
\end{exercise}

\begin{exercise}[$\star$]\label{exo:g2:using-a-computer:5}
Lina has finished her story, saved as \code{my story}. She wants to keep it as
it is and write a second version. Should she use \menu{File > Save} or
\menu{File > Save as}?
\end{exercise}

\begin{exercise}[$\star$]\label{exo:g2:using-a-computer:6}
What should Sami have done so that his castle drawing was still there the next
morning?
\end{exercise}

\begin{exercise}[$\star$]\label{exo:g2:using-a-computer:7}
In the drawing of the \omterm{def:g2:using-a-computer:desktop}{desktop}, which \omterm{def:g2:using-a-computer:window}{window} is in front? How can you bring the
other one to the front?
\end{exercise}

\begin{exercise}[$\star$]\label{exo:g2:using-a-computer:8}
Leo wants the title of his poem to be big and bold. Should he use a \omterm{def:g2:using-a-computer:word-processor}{text
editor} or a \omterm{def:g2:using-a-computer:word-processor}{word processor}? Why?
\end{exercise}

\begin{exercise}[$\star\star$]\label{exo:g2:using-a-computer:9}
Noor wrote two versions of a story: a short one on Monday and a longer one on
Wednesday. Give each version a good name.
\end{exercise}

\begin{exercise}[$\star\star$]\label{exo:g2:using-a-computer:10}
The \omterm{def:g2:using-a-computer:window}{window} of a tool stops answering: \omterm{def:g2:keyboard-and-mouse:click}{clicks} do nothing. What should you try
first? And then?
\end{exercise}

\begin{exercise}[$\star\star$]\label{exo:g2:using-a-computer:11}
Put these steps in order: \omterm{def:g2:keyboard-and-mouse:click}{double-click} on the name of your drawing; switch on
the \omterm{def:g1:what-is-a-computer:computer}{computer}; wait until the \omterm{def:g2:using-a-computer:desktop}{desktop} appears; open the place where you saved
the drawing; go on drawing; save.
\end{exercise}
